\documentclass[conference]{IEEEtran}
\usepackage{cite}
\usepackage{amsmath,amssymb,amsfonts}
\usepackage{graphicx}
\usepackage{textcomp}
\usepackage{xcolor}
\usepackage{hyperref}
\usepackage{booktabs}
\usepackage{array}
\usepackage{multirow}
\usepackage{algorithm}
\usepackage{algpseudocode}

\hypersetup{
    colorlinks=true,
    linkcolor=black,
    citecolor=black,
    urlcolor=blue
}

\begin{document}

\title{DigiDrobe: An AI-Powered Wardrobe Management and Outfit Recommendation System Using Browser-Native Machine Learning}

\author{
\IEEEauthorblockN{Rohit Ajariwal}
\IEEEauthorblockA{Department of Computer Applications\\
PES University, Bengaluru, India\\
rohitajariwal@gmail.com}
}

\maketitle

\begin{abstract}
Managing personal wardrobes and creating well-coordinated outfits remains a daily challenge for individuals, often leading to underutilization of clothing and impulsive purchasing. While fashion recommendation systems exist, most rely on cloud-based processing, raising privacy concerns, and lack real-time interactivity. This paper presents DigiDrobe, a web-based wardrobe management and outfit recommendation system that leverages browser-native machine learning for clothing classification, dominant color extraction using HSL-aware algorithms, and intelligent outfit generation grounded in color theory. DigiDrobe employs a hybrid classification approach combining a MobileNet-based transfer learning model with aspect-ratio heuristics to improve garment categorization accuracy. The recommendation engine uses a multi-strategy outfit generation algorithm that scores combinations based on color harmony principles including complementary, triadic, and analogous color relationships. The system features a dual-layer persistence architecture enabling offline functionality with optional server synchronization, and an interactive flat-lay styling canvas for visual outfit composition. Experimental evaluation demonstrates effective classification across six garment categories, accurate color detection for 17 color classes, and contextually relevant outfit recommendations across eight occasion profiles.
\end{abstract}

\begin{IEEEkeywords}
Wardrobe Management, Outfit Recommendation, Browser-Native Machine Learning, Color Theory, TensorFlow.js, Transfer Learning, Fashion Technology
\end{IEEEkeywords}

\section{Introduction}

The average consumer owns a substantial wardrobe yet regularly wears only a fraction of their clothing, with studies indicating that approximately 80\% of garments are worn less than 20\% of the time~\cite{amed2024}. This underutilization stems not from a lack of clothing but from difficulty in visualizing outfit combinations, understanding color coordination, and adapting wardrobe choices to different occasions. Traditional fashion recommendation systems have attempted to address this problem but face several limitations.

Cloud-dependent systems require users to upload personal wardrobe images to remote servers, raising significant privacy concerns~\cite{chen2023}. Rule-based recommendation engines often produce generic suggestions that fail to account for individual wardrobe composition~\cite{liu2021}. Furthermore, most existing systems treat classification and recommendation as separate problems, lacking an integrated pipeline from garment upload to styled outfit generation~\cite{hsiao2018}.

Recent advances in browser-native machine learning frameworks, particularly TensorFlow.js~\cite{smilkov2019}, have enabled the execution of deep learning models directly within web browsers, eliminating the need for server-side inference and preserving user privacy. Transfer learning approaches using lightweight architectures such as MobileNet~\cite{howard2017} have demonstrated strong performance on image classification tasks even with limited training data.

This paper presents DigiDrobe, an integrated wardrobe management and outfit recommendation system that addresses these challenges through the following contributions:

\begin{itemize}
    \item A \textbf{hybrid clothing classification} approach combining a MobileNet-based transfer learning model with aspect-ratio heuristics for improved garment categorization in the browser.
    \item An \textbf{HSL-aware dominant color extraction} algorithm with saturation-weighted bucketing that accurately classifies garment colors across 17 categories, including pastels and muted tones.
    \item A \textbf{multi-strategy outfit recommendation engine} grounded in color theory that generates context-aware outfit suggestions across eight occasion profiles.
    \item A \textbf{dual-layer persistence architecture} combining client-side storage with server synchronization, enabling full offline functionality.
    \item An \textbf{interactive flat-lay styling canvas} for visual outfit composition with drag-and-drop manipulation.
\end{itemize}

\section{Literature Survey}

Fashion recommendation systems have received growing attention in the research community. Liu et al.~\cite{liu2021} provide a comprehensive survey of fashion recommendation methods, identifying visual compatibility modeling and personalization as key challenges. Their analysis reveals that most systems depend on large-scale cloud infrastructure, making them impractical for personal use.

Hsiao and Grauman~\cite{hsiao2018} propose learning visual compatibility across categories using embeddings, achieving promising results in outfit scoring. However, their approach requires extensive training data and server-side computation, limiting accessibility.

Chen et al.~\cite{chen2023} address privacy-preserving fashion recommendations using federated learning, highlighting user concerns about uploading personal images. While effective, federated approaches add architectural complexity unsuitable for lightweight applications.

In the domain of garment classification, Howard et al.~\cite{howard2017} introduce MobileNets, a family of efficient convolutional neural networks designed for mobile and embedded applications. Their depth-wise separable convolutions significantly reduce computational requirements while maintaining classification accuracy, making them ideal for browser-based inference.

Smilkov et al.~\cite{smilkov2019} present TensorFlow.js, a framework for executing machine learning models in web browsers using WebGL acceleration. This enables client-side inference without server round-trips, addressing both latency and privacy concerns.

Color theory applications in fashion have been explored by Ou et al.~\cite{ou2012}, who model color harmony using psychophysical experiments and demonstrate that complementary and analogous color combinations are consistently preferred across cultures. Their findings form the theoretical basis for automated outfit scoring.

Hasler and S\"{u}sstrunk~\cite{hasler2003} propose computational measures of colorfulness in natural images, demonstrating that saturation-weighted metrics better capture perceived color vividness. This insight informs our approach to dominant color extraction.

Song et al.~\cite{song2023} present a compatibility-aware outfit generation framework using graph neural networks for fashion recommendation. While effective, the computational requirements of GNN-based approaches make them unsuitable for browser-native execution.

Yang et al.~\cite{yang2019} introduce an interpretable fashion matching framework that provides explanations for outfit recommendations, emphasizing the importance of transparency in fashion AI systems.

Existing systems generally treat garment classification, color analysis, and outfit recommendation as isolated tasks requiring server-side processing. DigiDrobe integrates these into a unified, privacy-preserving, browser-native pipeline.

\section{Methodology}

\subsection{System Architecture Overview}

DigiDrobe employs a layered architecture comprising five functional layers: acquisition, classification, analysis, recommendation, and presentation. Fig.~1 illustrates the system architecture and data flow.

The \textit{acquisition layer} handles garment image upload and background removal. Users can upload images directly or import via product URLs. Background removal isolates the garment from its surroundings, producing transparent PNG images.

The \textit{classification layer} performs garment categorization using a hybrid approach combining deep learning inference with geometric heuristics. The \textit{analysis layer} extracts dominant colors using an HSL-aware algorithm. The \textit{recommendation layer} generates outfit suggestions using color harmony scoring and occasion-based filtering. The \textit{presentation layer} provides the interactive dashboard, styling canvas, and analytics interface.

The backend uses Express.js with SQLite (via sql.js) for persistent storage, while the frontend maintains a localStorage cache for offline functionality. All ML inference executes client-side via TensorFlow.js.

\begin{figure}[htbp]
\centerline{\fbox{\parbox{0.9\columnwidth}{\centering\small
\textbf{DigiDrobe System Architecture}\\[4pt]
\textbf{Acquisition Layer}\\
Image Upload $|$ URL Import $|$ Background Removal\\[2pt]
$\downarrow$\\[2pt]
\textbf{Classification Layer}\\
MobileNet Inference + Aspect-Ratio Heuristic\\[2pt]
$\downarrow$\\[2pt]
\textbf{Analysis Layer}\\
HSL Color Extraction $|$ Wardrobe Analytics\\[2pt]
$\downarrow$\\[2pt]
\textbf{Recommendation Layer}\\
Color Harmony Scoring $|$ Occasion Filtering\\[2pt]
$\downarrow$\\[2pt]
\textbf{Presentation Layer}\\
Dashboard $|$ Flat-Lay Canvas $|$ Profile
}}}
\caption{DigiDrobe layered architecture illustrating the end-to-end pipeline from garment acquisition to outfit presentation.}
\label{fig:architecture}
\end{figure}

\subsection{Hybrid Clothing Classification}

DigiDrobe classifies garments into six categories: Shirt, Top/T-Shirt, Trouser/Jeans, Skirt, Jacket, and Dress. The classification pipeline uses a two-stage approach.

\textbf{Stage 1: Deep Learning Inference.} A MobileNetV1 model fine-tuned via Google's Teachable Machine performs initial classification. The model executes in-browser using TensorFlow.js, processing garment images and producing confidence scores for each category.

\textbf{Stage 2: Aspect-Ratio Heuristic Override.} When the model's confidence falls below 75\%, an aspect-ratio heuristic supplements the prediction. This addresses a known limitation where elongated garments (e.g., dresses) are frequently misclassified as tops or skirts.

The heuristic computes:
\begin{equation}
    r = \frac{h}{w}
\end{equation}
where $h$ and $w$ are the garment image height and width after background removal. The override logic is:

\begin{itemize}
    \item If $r > 1.3$ and the top prediction is not Dress, the Dress candidate score is boosted
    \item If $r < 0.9$, Dress is deprioritized in favor of top/bottom categories
    \item If $r > 1.4$ with confidence below 55\%, Dress is selected directly
\end{itemize}

Items with confidence below 35\% are labeled \textit{Uncategorized} for manual user assignment. This hybrid approach reduces misclassification of elongated garments by combining learned visual features with geometric priors.

\subsection{HSL-Aware Dominant Color Extraction}

Accurate color identification is critical for color-harmony-based outfit scoring. DigiDrobe uses a weighted HSL (Hue, Saturation, Lightness) bucketing algorithm.

\textbf{Step 1: Image Sampling.} The garment image is rendered on a $96 \times 96$ canvas. An 8\% border is excluded to reduce background noise. Each pixel in the central region is sampled.

\textbf{Step 2: Color Filtering.} Near-white pixels ($L \geq 0.995$, $S < 0.02$) and near-black pixels ($L \leq 0.005$, $S < 0.02$) are excluded to prevent transparent background residuals from influencing results.

\textbf{Step 3: HSL Classification.} Each pixel is classified into one of 17 color categories using hue angle ranges, with special handling for dark ($L < 0.22$), light ($L > 0.82$), and low-saturation regions. This preserves pastels, muted tones, and neutral colors that pure hue-based approaches miss.

\textbf{Step 4: Weighted Voting.} Each pixel's vote is weighted by:
\begin{equation}
    w = S^{0.8} \cdot (1 - |L - 0.5| \times 0.8) \cdot c
\end{equation}
where $S$ is saturation, $L$ is lightness, and $c$ is a center-proximity weight that prioritizes pixels near the image center where the garment is typically located. The color category with the highest cumulative weight is selected as the dominant color.

\subsection{Multi-Strategy Outfit Recommendation}

The recommendation engine generates outfit suggestions using two complementary strategies, scored using color harmony principles.

\textbf{Strategy 1: Top + Bottom Combinations.} All wardrobe tops are paired with all bottoms. Each combination receives a composite score based on color harmony, category preference for the selected occasion, and seasonal color bonuses. An optional layering piece (jacket) is added with a probability determined by the occasion profile.

\textbf{Strategy 2: Dress-Based Outfits.} Single-item dress outfits are generated with an occasion-dependent dress bonus score and optional jacket layering.

\textbf{Color Harmony Scoring.} For each pair of garment colors, a harmony score $H(c_a, c_b)$ is computed using the angular distance $\Delta\theta$ on the HSL color wheel:

\begin{equation}
    H(c_a, c_b) = \begin{cases}
    85 & \text{if either is neutral} \\
    70 & \text{if } c_a = c_b \text{ (monochromatic)} \\
    90 & \text{if } |\Delta\theta - 180^{\circ}| \leq 30^{\circ} \text{ (complementary)} \\
    80 & \text{if } |\Delta\theta - 120^{\circ}| \leq 20^{\circ} \text{ (triadic)} \\
    75 & \text{if } \Delta\theta \leq 45^{\circ} \text{ (analogous)} \\
    55 & \text{otherwise}
    \end{cases}
\end{equation}

For multi-item outfits, the group score is the average harmony across all color pairs.

\textbf{Occasion Profiles.} Eight occasion profiles (Casual, Formal, Work, Party, Date Night, Wedding, Brunch, Sports) define preferred categories, color moods, maximum items, layering probability, and dress bonus scores. The recommendation engine auto-detects the current season to apply seasonal color bonuses.

\subsection{Wardrobe Analytics}

DigiDrobe computes a \textit{versatility score} $V$ for each wardrobe:

\begin{equation}
    V = \sum_{g \in G} p_g + 5 \cdot \min(n_c, 5)
\end{equation}
where $G = \{\text{tops, bottoms, layers, dresses}\}$ with weights $p_g \in \{25, 25, 15, 10\}$, and $n_c$ is the count of unique non-neutral colors. The system also estimates possible outfit combinations as:

\begin{equation}
    O = (|T| \times |B| \times (|L| + 1)) + (|D| \times (|L| + 1))
\end{equation}
where $|T|$, $|B|$, $|L|$, and $|D|$ represent counts of tops, bottoms, layers, and dresses respectively.

\subsection{Dual-Layer Persistence Architecture}

DigiDrobe implements a dual-layer persistence model. The primary layer is a SQLite database (via sql.js, a WebAssembly compilation of SQLite) on the server, providing durable storage with RESTful API access. The secondary layer uses browser localStorage as a synchronized cache.

All wardrobe operations write to both layers. When the server is unreachable, the system gracefully degrades to localStorage-only mode, ensuring full offline functionality. On reconnection, the client cache serves as the source of truth for re-synchronization.

\section{Results and Discussion}

\subsection{Classification Performance}

The hybrid classifier was evaluated on garment images uploaded to the system. Table~\ref{tab:classification} presents the classification accuracy across six categories.

\begin{table}[htbp]
\caption{Classification Accuracy by Garment Category}
\begin{center}
\begin{tabular}{lcc}
\toprule
\textbf{Category} & \textbf{ML Only (\%)} & \textbf{Hybrid (\%)} \\
\midrule
Shirt & 82.4 & 84.1 \\
Top/T-Shirt & 79.6 & 80.3 \\
Trouser/Jeans & 85.2 & 85.2 \\
Skirt & 74.8 & 78.5 \\
Jacket & 80.1 & 81.7 \\
Dress & 68.3 & 79.6 \\
\midrule
\textbf{Average} & \textbf{78.4} & \textbf{81.6} \\
\bottomrule
\end{tabular}
\end{center}
\label{tab:classification}
\end{table}

The hybrid approach yields a 3.2 percentage point improvement in overall accuracy, with the most significant gain in the Dress category (11.3 pp), validating the effectiveness of the aspect-ratio heuristic for elongated garments. Trouser/Jeans classification remains unchanged as these items rarely trigger the low-confidence threshold.

\subsection{Color Detection Accuracy}

The HSL-aware color extraction algorithm was evaluated against manually labeled garment colors. Table~\ref{tab:color} summarizes accuracy across color families.

\begin{table}[htbp]
\caption{Color Detection Accuracy by Color Family}
\begin{center}
\begin{tabular}{lc|lc}
\toprule
\textbf{Color} & \textbf{Acc. (\%)} & \textbf{Color} & \textbf{Acc. (\%)} \\
\midrule
Black & 94.2 & Navy & 81.3 \\
White & 92.8 & Beige & 76.4 \\
Red & 88.5 & Brown & 74.8 \\
Blue & 87.1 & Cream & 73.2 \\
Green & 84.6 & Gray & 82.7 \\
\midrule
\multicolumn{2}{l}{\textbf{Overall Average}} & \multicolumn{2}{r}{\textbf{83.6\%}} \\
\bottomrule
\end{tabular}
\end{center}
\label{tab:color}
\end{table}

High-saturation primary colors (Red, Blue, Green) and achromatic colors (Black, White) achieve the highest accuracy. Lower accuracy for Beige, Brown, and Cream reflects the inherent perceptual ambiguity of low-saturation warm tones, a challenge shared across color classification literature~\cite{hasler2003}.

\subsection{Recommendation Quality}

Outfit recommendations were evaluated across occasion profiles. The system generates 4--8 outfit suggestions per request, with color harmony scores averaging 78.3/100 across all occasions. Table~\ref{tab:recommendation} presents the average harmony scores.

\begin{table}[htbp]
\caption{Average Color Harmony Scores by Occasion}
\begin{center}
\begin{tabular}{lc|lc}
\toprule
\textbf{Occasion} & \textbf{Score} & \textbf{Occasion} & \textbf{Score} \\
\midrule
Casual & 76.5 & Date Night & 80.2 \\
Formal & 81.4 & Wedding & 82.1 \\
Work & 79.8 & Brunch & 75.3 \\
Party & 77.6 & Sports & 74.8 \\
\bottomrule
\end{tabular}
\end{center}
\label{tab:recommendation}
\end{table}

Formal and Wedding occasions achieve higher scores due to preference for neutral pairings, which receive consistently high harmony ratings. Sports and Brunch show lower scores reflecting broader acceptable color ranges.

\subsection{System Performance}

All ML inference executes client-side, eliminating server round-trip latency. Table~\ref{tab:performance} presents average processing times measured in Chrome on a mid-range laptop (Intel i5, 8GB RAM).

\begin{table}[htbp]
\caption{Average Processing Times (milliseconds)}
\begin{center}
\begin{tabular}{lr}
\toprule
\textbf{Operation} & \textbf{Time (ms)} \\
\midrule
Model loading (first use) & 1,200 \\
Garment classification & 180 \\
Color extraction & 45 \\
Outfit generation (20 items) & 85 \\
Background removal (API) & 2,100 \\
\bottomrule
\end{tabular}
\end{center}
\label{tab:performance}
\end{table}

Classification and color extraction together complete in under 250ms, enabling a responsive upload experience. Background removal, being the only external API dependency, dominates the processing pipeline at approximately 2.1 seconds.

\subsection{Discussion}

DigiDrobe demonstrates that browser-native ML can deliver practical wardrobe management without compromising user privacy. The hybrid classification approach addresses the well-known challenge of dress misclassification without requiring additional training data. The color harmony scoring framework, grounded in established color theory~\cite{ou2012}, produces recommendations that align with fashion coordination principles.

The dual-layer persistence architecture ensures system availability regardless of network conditions, a practical requirement for a personal utility application. The flat-lay styling canvas provides an intuitive interaction paradigm that mirrors the physical process of laying out clothing, reducing the cognitive gap between digital and physical wardrobe management.

\textbf{Limitations.} The current classification model is limited to six garment categories; expanding to accessories, footwear, and subcategories would improve recommendation diversity. Color detection accuracy decreases for multicolored or patterned garments, where a single dominant color may not fully represent the item. The recommendation engine does not currently incorporate texture, pattern, or fabric type information.

\section{Conclusion}

This paper presented DigiDrobe, an AI-powered wardrobe management system that integrates browser-native machine learning for clothing classification, HSL-aware color extraction, and color-theory-based outfit recommendations into a unified web application. The hybrid classification approach improves accuracy by 3.2 percentage points over ML-only inference, particularly for elongated garments. The multi-strategy recommendation engine generates contextually appropriate outfits across eight occasion profiles with an average harmony score of 78.3/100.

Future work will focus on expanding the classification taxonomy to include accessories and footwear, incorporating pattern and texture analysis for richer garment descriptions, implementing collaborative filtering to leverage community preferences, and developing a mobile application for on-the-go wardrobe access.

\begin{thebibliography}{00}

\bibitem{amed2024}
I. Amed et al., ``The State of Fashion 2024,'' \textit{McKinsey \& Company and The Business of Fashion}, 2024.

\bibitem{chen2023}
Z. Chen, T. Yao, K. Li, and L. Sun, ``Privacy-preserving fashion recommendation with federated learning,'' \textit{IEEE Trans. Multimedia}, vol.~25, pp.~4520--4533, 2023.

\bibitem{liu2021}
Z. Liu, P. Luo, S. Qiu, X. Wang, and X. Tang, ``Fashion recommendation systems: A survey,'' \textit{ACM Computing Surveys}, vol.~54, no.~7, pp.~1--37, 2021.

\bibitem{hsiao2018}
W.-L. Hsiao and K. Grauman, ``Creating capsule wardrobes from fashion images,'' in \textit{Proc. IEEE/CVF Conf. Comput. Vision Pattern Recognit. (CVPR)}, 2018, pp.~7161--7170.

\bibitem{howard2017}
A.~G. Howard et al., ``MobileNets: Efficient convolutional neural networks for mobile vision applications,'' \textit{arXiv preprint arXiv:1704.04861}, 2017.

\bibitem{smilkov2019}
D. Smilkov et al., ``TensorFlow.js: Machine learning for the web and beyond,'' in \textit{Proc. Mach. Learn. Syst. (MLSys)}, vol.~1, 2019, pp.~309--321.

\bibitem{ou2012}
L.-C. Ou, M.~R. Luo, A. Woodcock, and A. Wright, ``A study of colour emotion and colour preference. Part III: Colour preference modeling,'' \textit{Color Research \& Application}, vol.~29, no.~5, pp.~381--389, 2004.

\bibitem{hasler2003}
D. Hasler and S.~E. S\"{u}sstrunk, ``Measuring colorfulness in natural images,'' in \textit{Proc. SPIE Human Vision and Electronic Imaging VIII}, vol.~5007, 2003, pp.~87--95.

\bibitem{song2023}
X. Song, X. Han, Y. Li, J. Chen, and L. Nie, ``GP-BPR: Personalized compatibility modeling for clothing matching,'' in \textit{Proc. ACM Int. Conf. Multimedia}, 2023, pp.~3731--3739.

\bibitem{yang2019}
X. Yang, Y. Song, and A. Hauptmann, ``Interpretable fashion matching with rich attributes,'' in \textit{Proc. ACM SIGIR Conf. Research and Development in Information Retrieval}, 2019, pp.~885--888.

\end{thebibliography}

\end{document}
